\documentclass[11pt]{article}
\usepackage[margin=1in]{geometry}
\usepackage[T1]{fontenc}
\usepackage{longtable}
\usepackage{array}
\begin{document}
\section*{FLCR-02 - The LCR Carrier}

**Series:** Formalizing LCR, Unifying Standard Models
**Artifact:** formal paper source
**Status:** first-pass rich rewrite; requires final citation and build pass.
**Claim posture:** maximal local claim posture with explicit lane boundaries.

\subsection*{Abstract}

This paper defines the LCR carrier as the canonical radius-1 local readout with a distinguished center and two ordered boundary slots. The carrier is finite, inspectable, and minimal: any smaller carrier collapses either the center or one boundary address. The paper proves the binary eight-state surface, reversal action, shell grading, and shell-2 doublet-plus-pivot structure used by later papers.

\subsection*{Keywords}

LCR carrier; radius-1 cellular automaton; C2 action; shell grading; minimal arity.

\subsection*{External Reader Orientation}

An outside reviewer should read this paper as a formal-system construction paper. Its local object is **The LCR Carrier**. The paper's immediate contribution is: Defines the L/C/R carrier as the first active state surface.

The nearest external anchors are finite-state reasoning, typed interfaces, proof-carrying records, conservative extension, and ordinary mathematical citation discipline. LCR terms are therefore internal labels for the construction unless a row explicitly imports an external theorem, citation, dataset, or calibration target.

It does not ask the reader to accept any Standard Model or literal-physics identification at this layer. Such mappings are deferred to the translation papers and must carry their own evidence.

\subsection*{Reviewer Compression}

**Core object.** The LCR Carrier.

**Primary result.** This paper contributes the local formal object and its safe downstream-use rule.

**Primary non-result.** This paper does not claim direct identity with Standard Model or physical objects.

**Review strategy.** Evaluate the formal claim rows first, then inspect the receipt/citation/calibration rows that support each claim. Appendices provide routing and implementation context; they should not be read as stronger than their lane permits.

\subsection*{Peer Reviewer Assessment}

Reviewer assumption: the reviewer has been briefed on the FLCR structure, claim lanes, CAM/crystal routing, forge/action surfaces, and the distinction between internal formal results and external physical calibration.

**Recommendation.** major revision, technically promising.

**Review score vector.** orientation=2, claim\_granularity=2, evidence\_visibility=2, falsifier\_visibility=2, journal\_apparatus=1, base\_layer\_boundary=2.

**Formal claim rows reviewed.** 3.

\subsubsection*{Reviewer Strength Finding}

The manuscript is now readable as a bounded formal-system paper: it identifies its local object, separates internal construction from physical interpretation, and exposes claim lanes for review.

\subsubsection*{Major Revision Requests}

\noindent{}-- Convert the strongest proof sketch into numbered definitions, lemmas, and propositions with explicit dependencies.\\
\noindent{}-- Replace placeholder source classes with final citation keys, receipt hashes, or dataset identifiers wherever possible.\\
\noindent{}-- Move repeated pass-derived material into a compact evidence appendix and keep the main body focused on the paper-local argument.\\
\noindent{}-- Add at least one worked example showing the local object, legal operation, receipt, residue, and downstream import rule.\\

\subsubsection*{Minor Revision Requests}

\noindent{}-- Normalize theorem capitalization and punctuation.\\
\noindent{}-- Add final figure/table captions where the text currently describes diagrams schematically.\\
\noindent{}-- Ensure every acronym and local term appears in the paper-local glossary.\\

\subsubsection*{Acceptance Blockers}

\noindent{}-- A reviewer can accept the formal contribution without accepting later physics claims, but only if final citations and receipt identifiers are attached.\\
\noindent{}-- The proof sketch must be promoted into a formal proof or explicitly labeled as a construction argument.\\


\subsection*{1. Contribution And Scope}

\noindent{}-- Defines LCR as the primitive local address surface for the series.\\
\noindent{}-- Proves the three-slot arity lower bound for center-plus-two-boundary data.\\
\noindent{}-- Registers the binary shell profile 1,3,3,1 and the reversal action.\\
\noindent{}-- Separates internal spinor-like local semantics from physical spin transport.\\

This paper belongs to the LCR-native construction band. Physics and Standard Model language, when it appears, is only a deferred mapping cue, analogy, or explicitly bounded calibration target. The paper's base object must remain stable enough for later papers to cite without importing a stronger physical claim by implication.

\subsubsection*{Series Posture Inherited From FLCR-01}

This paper inherits the FLCR-01 doctrine. Guardrails are automation controls:
they govern AI agents, validators, build scripts, generated prose, and claim
promotion workflows. They are not restrictions on human creativity or
hypothesis formation.

The paper should therefore name the strongest intended reading of `The LCR Carrier`,
display the parts already proved at this layer, and route the remaining
distance as explicit open obligations. Those obligations are channels from the
current prestate into later exploration.

The analog/workbook layer is also an access kit. It should preserve the local
state in forms that can eventually be replayed without electricity, internet,
computers, or paper. Later papers may work toward literal physics even where
this local paper only supplies the finite or LCR-native carrier.

\subsection*{2. Source Routing}

Primary routed sources:

\noindent{}-- `01\_LCR\_Chain\_Carrier.md`\\
\noindent{}-- `supplements/01\_INTERNAL\_REASONING\_SUPPLEMENT.md`\\

Cross-corpus evidence classes:

\noindent{}-- `CAM\_CLAIM\_100\_SCORE\_AUDIT.md`\\
\noindent{}-- `NSIT\_TOOL\_CLOSURE\_MATRIX.md`\\
\noindent{}-- `NSIT\_REASONED\_CLOSURE\_PASS.md`\\
\noindent{}-- `PAPER\_REWORKS\_CRYSTAL\_PROJECTION.json`\\
\noindent{}-- standards, glue, window, and node databases where applicable\\
\noindent{}-- notebook-derived review prompts and orbital source manifests\\

\subsection*{3. Definitions}

\noindent{}-- **LCR state.** A triple (L, C, R) with left boundary, center, and right boundary coordinates.\\
\noindent{}-- **Binary chart.** The finite carrier \{0,1\}\textasciicircum{}3 containing exactly eight states.\\
\noindent{}-- **Shell grade.** The Hamming weight of the binary state, giving multiplicities 1,3,3,1.\\
\noindent{}-- **Reversal.** The C2 action rho(L,C,R) = (R,C,L).\\

\subsection*{4. Formal Claims}

\noindent\texttt{| Claim | Lane | Statement |}\\
\noindent\texttt{|-------|------|-----------|}\\
\noindent\texttt{| Theorem 2.1 | `receipt\_bound\_internal\_result` | The LCR carrier is the minimal address-preserving carrier for one center and two distinguishable boundary slots. |}\\
\noindent\texttt{| Theorem 2.2 | `receipt\_bound\_internal\_result` | In the binary chart, reversal is an involution and partitions the carrier into fixed states and reversal pairs. |}\\
\noindent\texttt{| Proposition 2.3 | `standard\_theorem\_citation\_bound\_result` | The shell profile of \{0,1\}\textasciicircum{}3 is the binomial profile 1,3,3,1. |}\\

\subsection*{Reviewer Claim Ledger}

This ledger expands the formal-claims table into review actions. It is intended to make proof granularity explicit: what is being claimed, what evidence type can support it, what boundary remains, and what the next review action is.

\noindent\texttt{| Formal row | Type | Claim in review terms | Evidence required | Boundary | Next review action |}\\
\noindent\texttt{|---|---|---|---|---|---|}\\
\noindent\texttt{| Theorem 2.1 | finite/executable result | The LCR carrier is the minimal address-preserving carrier for one center and two distinguishable boundary slots | receipt, validator, executable pass, generated artifact, or fin}\\
\noindent\texttt{| Theorem 2.2 | finite/executable result | In the binary chart, reversal is an involution and partitions the carrier into fixed states and reversal pairs | receipt, validator, executable pass, generated artifact, or fini}\\
\noindent\texttt{| Proposition 2.3 | formal construction | The shell profile of \{0,1\}\textasciicircum{}3 is the binomial profile 1,3,3,1 | named external theorem, definition layer, citation target, or imported standard formalism | imports the external ob}\\

\subsection*{Claim Granularity And Review Table}

The table below separates the claim types so the paper is reviewable without accepting the whole FLCR vocabulary at once.

\noindent\texttt{| Claim type | What this paper may claim | Acceptance test | What is not claimed by that row |}\\
\noindent\texttt{|---|---|---|---|}\\
\noindent\texttt{| Formal-system result | `FLCR-02` defines or uses **The LCR Carrier** as a typed FLCR object with declared inputs, operations, outputs, and residue. | Definitions, formal claims, construction steps, and downstream-use r}\\
\noindent\texttt{| Computational or receipt-bound result | Enumerations, TarPit runs, CAM/crystal routes, verifier rows, and generated artifacts are claims about the stated finite or executable domain. | A named receipt, validator, manif}\\
\noindent\texttt{| Imported theorem or external definition | Imported mathematics remains an external theorem/citation target; this paper may use it only through declared mappings and conservative-extension discipline. | The source theor}\\
\noindent\texttt{| Interpretive bridge or analogy | Analog, workbook, visual, or translation language may explain why the construction is useful. | The analogy preserves the relevant structure and does not promote itself into a theorem. }\\
\noindent\texttt{| Physical or calibration-facing claim | Physical or Standard Model identity is not claimed here unless the paper contains an explicit calibration row; the default status is deferred mapping. | A dataset, citation, calib}\\
\noindent\texttt{| Open obligation or falsifier | A missing derivation, adapter, receipt, dataset, or failed comparator is a named research channel. | The next binding action and the reason the stronger claim is not closed are stated. | }\\

\subsection*{5. Paper-Specific Development}

1. Represent each local observation as (L,C,R), preserving left boundary, center, and right boundary.
2. Enumerate the binary chart \{0,1\}\textasciicircum{}3 and compute shell grade by Hamming weight.
3. Apply reversal rho(L,C,R)=(R,C,L) and record fixed points and two-state orbits.
4. Use the arity lower bound to show why fewer than three slots destroys either center identity or boundary distinction.

\subsubsection*{Proof Sketch}

The finite proof is direct enumeration plus a minimality argument. There are exactly eight binary triples. Reversal applied twice returns the original triple, so it is a C2 action. Any two-slot representation has only two addresses and therefore must identify two of the three roles \{left, center, right\}; this violates the carrier contract.

\subsection*{6. Construction}

The construction is intentionally staged. First, the paper names the finite or
typed object that can be inspected directly. Second, it declares the admissible
operations over that object. Third, it records the receipt or source class that
allows another paper to consume the result. Fourth, it names the residue that
must not be silently promoted.

For this paper, the operative object is `The LCR Carrier`. Its safe consumption
rule is:

\begin{verbatim}
source object -> declared operation -> receipt/lane -> downstream import
\end{verbatim}

The unsafe consumption rule is:

\begin{verbatim}
source phrase -> downstream identity claim
\end{verbatim}

That second pattern is explicitly rejected by FLCR-01 and must be rewritten as
a claim-lane promotion with evidence.

\subsection*{Recent External Quantum Evidence Intake}

These rows add newly routed external physics papers as citation-bound or calibration-bound evidence. They are not CAM receipts unless a later tool run reconstructs their signatures.

\noindent\texttt{| Source | Lane | How It Helps This Paper | Boundary |}\\
\noindent\texttt{|---|---|---|---|}\\
\noindent\texttt{| EXT-RQM-CONSISTENT-2026: Quantum mechanics based on real numbers: A consistent description; DOI `10.1103/4k13-sdjh`; arXiv `2503.17307` | `standard\_theorem\_citation\_bound\_result` | Supports the carrier-layer distinctio}\\

\subsection*{7. Evidence And Receipts}

\noindent\texttt{| Evidence source | What it contributes |}\\
\noindent\texttt{|-----------------|---------------------|}\\
\noindent\texttt{| Paper 01 legacy body | Carrier enumeration, shell-2 doublet, local O8/double-cover route. |}\\
\noindent\texttt{| Internal supplement 01 | Radius-1 CA grounding, arity lower bound, C2 action, binomial shell grading. |}\\
\noindent\texttt{| NSIT tool matrix | T5-T7 plus shell-2 doublet receipt and O8 spinor double-cover closure route. |}\\
\noindent\texttt{| CAM Claim 100 Score Audit | Paper 01 scored 94 as internally closed, with physical spinor transport deferred. |}\\

\subsection*{Appendix C. Implementation Intake}

This addendum binds implementation work that was not fully represented in the earlier paper routing. The rows below are not pasted source text; they are evidence-lane imports that change what the paper can claim now.

\noindent\texttt{| Implementation source | Lane | Claim effect | Boundary |}\\
\noindent\texttt{|---|---|---|---|}\\
\noindent\texttt{| IMPL-LCR-TYPED-KERNEL | `normal\_form\_result` | Promotes L/C/R from notation to operational lane contract: LAdapter, CKernel, RChannel, strict policy gates, and lane classification. | This closes operational admission a}\\
\noindent\texttt{| IMPL-R30-LATTICE-THEOREM-REGISTRY | `standard\_theorem\_citation\_bound\_result` | Adds executable theorem registry rows for octonions, J3(O), chart-to-J3(O), exact n=3 SU(3) Weyl closure, Rule-30 8x8 transition, Niemeier }\\

The immediate paper effect is stronger claim posture where these implementation rows satisfy the relevant lane. Remaining caution is reserved for specific claims that still lack an adapter, comparator, calibration target, or rerun receipt.

\subsection*{Appendix D. Resolved-State Closure Pass}

This pass removes false restrictions on the paper's claim posture. A row is no longer called open merely because a stronger future claim exists. The satisfied lane is closed now; only the stronger claim that lacks its required adapter, comparator, calibration, or falsifier remains as residue.

\subsubsection*{Closed Now}

\noindent\texttt{| Row | Lane | Resolved state | Remaining boundary |}\\
\noindent\texttt{|---|---|---|---|}\\
\noindent\texttt{| paper lift enumeration | `receipt\_bound\_internal\_result` | 8/8 LCR states succeeded; error walls=0; avg TarPit mass=0.526710. | This closes the paper-local finite lift readout, not every stronger physical interpretatio}\\
\noindent\texttt{| carrier enumeration | `normal\_form\_result` | The LCR carrier action is closed as an addressable finite-state surface for its declared domain. | The family closure is a structural lane; measured physics still requires d}\\
\noindent\texttt{| decade-1 tower | `receipt\_bound\_internal\_result` | The decade tower is resolved with avg TarPit mass=0.510370 and mass sum=40.829567. | The decade total is an internal tower metric, not a standalone universal physical }\\
\noindent\texttt{| family-02 cross-lift comparison | `cam\_crystal\_reapplication\_result` | Family tower binds FLCR-02, FLCR-12, FLCR-22, FLCR-32; avg TarPit mass=0.511700; error walls=0. | Cross-lift agreement strengthens routing and comp}\\
\noindent\texttt{| finite LCR carrier | `receipt\_bound\_internal\_result` | The 2x2x2 carrier is closed as the finite local object consumed by later papers. | Physical interpretation of a carrier requires adapter rows. |}\\

\subsubsection*{Claim Posture After This Pass}

`FLCR-02` should now be read as a resolved-state contribution for `The LCR Carrier` at its declared lane. The paper may state the strongest claim supported by these rows directly. It should reserve caution only for a specifically named stronger claim whose required evidence is absent.

\subsection*{8. Dependencies And Downstream Use}

This paper may be imported by later FLCR papers only through the claim lanes
above. The default downstream consumers are:

\noindent{}-- `FLCR-11` through `FLCR-18` for window, receipt, admission, CA, algebraic,\\
  curvature, residue, and exceptional machinery.
\noindent{}-- `FLCR-28` through `FLCR-30` for CAM/crystal, corpus ribbon, and universal\\
  normal-form intake.
\noindent{}-- `FLCR-31` through `FLCR-40` only after the LCR-native result has been cited\\
  and the translation claim has its own evidence lane.

\subsection*{9. Limitations And Falsifiers}

\noindent{}-- The shell-2 doublet is an internal interface, not a measured particle claim.\\
\noindent{}-- SU(2)-style language may be used as standard analogy only where explicitly lane-tagged.\\
\noindent{}-- Later observer or physics papers must attach calibration or observer receipts before consuming this as physical spin.\\

A falsifier for this paper is any example that satisfies the declared input
conditions while violating the stated construction, receipt, or lane boundary.
An interpretation that merely wants a stronger downstream conclusion is not a
falsifier; it is an obligation routed to a later paper.

\subsection*{10. Publication Readiness Tasks}

\noindent{}-- Validator identities and receipt hashes are recorded in the manifest or receipt appendix.\\
\noindent{}-- External theorems require final citation entries before journal submission.\\
\noindent{}-- Every theorem, proposition, and protocol must retain a claim lane and evidence boundary.\\
\noindent{}-- The Markdown, LaTeX, and PDF forms are generated from the same canonical source.\\

\subsection*{Dual Positioning: Story And Formal Carrier}

This paper must be read in two synchronized ways. Sequentially, it explains why
the next state exists in the corpus story. Formally, it binds that state to
accepted mathematics, IRL formalism, code receipts, validators, CAM/crystal
lookups, and claim-lane boundaries.

Assigned original ribbon role(s): `01`/enumeration\_action.

\noindent\texttt{| Original slot | Family | Lift | Role | Proof form |}\\
\noindent\texttt{|---|---|---:|---|---|}\\
\noindent\texttt{| `01` | enumeration\_action | 0 | order-1 slot-1: start the active one-dimensional action / enumerate the center state | enumeration proof and identity preservation |}\\

The formal obligation is to state the strongest valid claim available for this
paper without letting interpretation silently change the addressed object. Any
author interpretation belongs beside the formalism as a declared relabeling,
bridge, or analog, and must be accompanied by the evidence lane that makes the
claim consumable by later papers.

\subsection*{State-Bound Proof Contract}

This paper receives: state emitted by prior slot 00 and reopened at original slot 01 (LCR Chain Carrier).

It must produce: formal result of original slot 01 plus the same-family lift path toward slot 11.

Semantic transition: select the active center and enumerate the addressable state space at the current lift.

Accepted formalism targets to bind in the journal rewrite:

\noindent{}-- finite set enumeration\\
\noindent{}-- ordered tuples and projections\\
\noindent{}-- group actions on finite carriers\\
\noindent{}-- minimality or lower-bound arguments\\

\noindent\texttt{| Slot | Family | Lift | Produced proof form |}\\
\noindent\texttt{|---|---|---:|---|}\\
\noindent\texttt{| `01` | enumeration\_action | 0 | enumeration proof and identity preservation |}\\

The formal carrier uses these accepted tools while preserving interpretive
claims as explicit relabelings, correspondences, or bridges. Claim promotion is admissible only when a receipt, standard citation,
CAM/crystal reapplication, normal-form result, calibration datum, or falsifier
boundary is attached.

\subsection*{Provenance And Crystal Evidence Integration}

This section integrates prior work, CAM/crystal projections, NSIT tooling,
standards-window evidence, and routed supplements as provenance-bearing
evidence. Source material is not treated as manuscript text; it is bound into
the paper only through claim lanes, receipts, validators, normal forms,
calibration rows, or explicit falsifier boundaries.

\subsubsection*{Expansion Rule For This Slot}

Use past work to increase the state inventory: enumerate the local carrier, admission candidates, or applied reader objects before interpretation is added.

\subsubsection*{Primary Evidence Inputs}

\noindent\texttt{| Source | Placement reason |}\\
\noindent\texttt{|---|---|}\\
\noindent\texttt{| `01\_LCR\_Chain\_Carrier.md` | primary assigned source for the paper's formal spine |}\\
\noindent\texttt{| `supplements/01\_INTERNAL\_REASONING\_SUPPLEMENT.md` | paper-local reasoning supplement contributing to definitions, proof sketch, and limitations |}\\

\subsubsection*{Secondary And Orbital Evidence Inputs}

\noindent\texttt{| Source | Placement reason |}\\
\noindent\texttt{|---|---|}\\
\noindent\texttt{| `supplements/RECEIPT\_VERIFIER\_CATALOG.md` | cross-cutting supplement; paper-relevant rows, receipts, and guard language are bound through evidence lanes |}\\
\noindent\texttt{| `supplements/INTERNAL\_REASONING\_PYRAMID\_INDEX.md` | cross-cutting supplement; paper-relevant rows, receipts, and guard language are bound through evidence lanes |}\\
\noindent\texttt{| `supplements/INTERNAL\_REASONING\_5P\_WINDOW\_SUPPLEMENT.md` | cross-cutting supplement; paper-relevant rows, receipts, and guard language are bound through evidence lanes |}\\
\noindent\texttt{| `supplements/INTERNAL\_REASONING\_7P\_WINDOW\_SUPPLEMENT.md` | cross-cutting supplement; paper-relevant rows, receipts, and guard language are bound through evidence lanes |}\\
\noindent\texttt{| `supplements/INTERNAL\_REASONING\_9P\_WINDOW\_SUPPLEMENT.md` | cross-cutting supplement; paper-relevant rows, receipts, and guard language are bound through evidence lanes |}\\
\noindent\texttt{| `CAM\_CLAIM\_100\_SCORE\_AUDIT.md` | series-wide audit/score/closure evidence; import paper-specific rows and leave global context outside the body |}\\
\noindent\texttt{| `NSIT\_TOOL\_CLOSURE\_MATRIX.md` | series-wide audit/score/closure evidence; import paper-specific rows and leave global context outside the body |}\\
\noindent\texttt{| `NSIT\_REASONED\_CLOSURE\_PASS.md` | series-wide audit/score/closure evidence; import paper-specific rows and leave global context outside the body |}\\
\noindent\texttt{| `ORBITAL\_DATA\_CROSS\_POPULATION\_AUDIT.md` | series-wide audit/score/closure evidence; import paper-specific rows and leave global context outside the body |}\\
\noindent\texttt{| `AGENT\_CRYSTAL\_WORK\_HARVEST.md` | series-wide audit/score/closure evidence; import paper-specific rows and leave global context outside the body |}\\
\noindent\texttt{| `PAPER\_REWORKS\_CRYSTAL\_PROJECTION.json` | series-wide audit/score/closure evidence; import paper-specific rows and leave global context outside the body |}\\
\noindent\texttt{| `NEW\_PAPER\_NEEDS.md` | route relevant obligations into this paper's burden table: NP-04 F4 Encoder Universality And Exceptional Route Conditions |}\\
\noindent\texttt{| Additional source routes | Complete routing is maintained in the source-placement appendix |}\\

\subsubsection*{Crystal And Standards Evidence To Bind}

\noindent{}-- Paper Reworks crystal projection: 33 paper markdown files, 9 supplements, 5 queues, 6 lattice-forge unification artifacts, 140 faces, 140 vignettes, and 280 CAM nodes.\\
\noindent{}-- Standards-window intake: 192/192 corpus conformance verdicts PASS; 140/142 obligations have candidate routes; 2/4/8/16/32 window lattice is available for cross-paper reads.\\
\noindent{}-- Agent/crystal harvest: agent-generated memory, runtime standards starter, current runtime build, repo-harvest CAM, notebook-derived bundles, and downloaded crystal payloads are orbital evidence surfaces.\\
\noindent{}-- NSIT inventory baseline: 220 validators, 1786 receipt/data artifacts, 1137 formal-paper-like artifacts, 114 supplements, and 86 memory/CAM artifacts.\\

\subsubsection*{Paper-Specific CAM/Score Rows}

\noindent\texttt{| 01 | 94 | internally closed | LCR minimality, shell-2 doublet, local O8 double cover, Paper 03 carry-forward | physical spinor transport remains NP-10/physics bridge |}\\

\subsubsection*{Paper-Specific NSIT Closure Rows}

\noindent\texttt{| T5-T7 plus shell-2 doublet | batch\_tool\_unison\_closure | `f4\_action.py`, `core.py` | Papers 03 and 01 | `su3\_closure\_T5\_T7\_receipt.json` 10/10; `bijective\_shell2\_doublet\_receipt.json` 7/7 | SU(3)/8x8 closure and single}\\
\noindent\texttt{| O8 spinor SU(2) double cover | batch\_tool\_unison\_closure | `oloid\_kinematic` frame inversion semantics | Paper 01 | `o8\_spinor\_double\_cover\_closed\_receipt.json` 6/6 | F\textasciicircum{}2 = -1 at 2pi; F\textasciicircum{}4 = +1 at 4pi. |}\\

\subsubsection*{Source Integration Summary}

The legacy source and internal reasoning supplement are treated as source
evidence rather than manuscript prose. Their contributions enter this paper as
claim-lane entries, receipt or validator references, finite-domain statements,
and limitation boundaries. Speculative or cross-domain interpretations remain
separate from closed local results until an explicit adapter, citation,
normal-form derivation, calibration row, or falsifier is attached.
\subsection*{Claim Binding Queue}

This queue records the evidence discipline for claim strengthening. It separates claims that already have support from residues that remain in falsifier or open-obligation lanes.

\noindent\texttt{| Queue | Lane | Promote now | Bind next | Residue |}\\
\noindent\texttt{|---|---|---|---|---|}\\
\noindent\texttt{| Spinor/SU(2) Double-Cover Local Semantics | `receipt\_bound\_internal\_result` | Promote local 2pi sign-flip / 4pi return semantics where the O8/frame-inversion receipt is attached. | Bind the O8 receipt and cite standard}\\

\subsubsection*{Binding Discipline}

Each queue item is represented as a delta:

\begin{verbatim}
local claim at paper time
-> attached later source/tool/receipt/theorem
-> revised representation
-> invariant boundary and remaining residue
\end{verbatim}

This preserves the sequential story while allowing later tools, crystals, and
standard formalisms to return through the proper lane.

\subsection*{Claim Delta Ledger}

This ledger applies the paper's claim-binding queues as explicit back-application
deltas. It keeps the sequential paper story intact while allowing later
receipts, standard formalisms, CAM/crystal routes, and validators to sharpen the
claim.

\noindent\texttt{| Delta | Local claim at paper time | Later evidence | Revised representation | Boundary kept |}\\
\noindent\texttt{|---|---|---|---|---|}\\
\noindent\texttt{| Spinor Double Cover Local | This paper may claim local frame-inversion or spinor-style behavior only for its finite/internal carrier object. | O8 frame-inversion receipt plus standard SU(2)->SO(3) double-cover semantic}\\

The revised representation records the strongest claim supported by the current evidence package. Promotion into theorem or proposition language occurs only when the named evidence is attached in the manifest or workbook replay.

\subsection*{Source-Backed Evidence Summary}

Routed archive and mirror cards are treated as provenance summaries rather than
body text. They identify candidate receipts, validators, theorem registries, or
supplemental source routes. A card strengthens this paper only when its
contribution is bound to a claim lane, evidence object, and boundary.

\noindent\texttt{| Evidence card | Score | Publication contribution | Boundary |}\\
\noindent\texttt{|---|---:|---|---|}\\
\noindent\texttt{| CQE-paper-10: Paper 10 - T10 Master Receipt | 81 | Paper 10 proves the first stack-level receipt theorem in the CQECMPLX suite. Its object is not a new physical mechanism. Its object is the proof that Paper 00 and Pape}\\
\noindent\texttt{| FORMAL: Paper 24 — C-Form: KnightForge / N-Dimensional Chess Automata Gluon | 77 | **C = the chess automata Gluon** — the L-conjugate CA Gluon that generalizes knight moves to N-dimensional board operators. In the latt}\\
\noindent\texttt{| FINAL\_FORMAL\_PAPER: Complete Formal Claims: The Folded Strand | 75 | We present the complete closed-form claim set of the CQE\_CMPLX corpus: - 33 papers = 33 folding operations on a self-complementary strand - 144 tools}\\
\noindent\texttt{| SUMMARY-V-32-Theorems-Registry: Summary Paper V — The 32 Theorems in One Table: Closed-Form Registry | 75 | This paper is the **complete closed-form registry of all 32 theorems** in the CQE\_CMPLX corpus. Each theorem i}\\
\noindent\texttt{| CQE-paper-10.25: Paper 10.25 - Toolkit for the T10 Master Receipt | 73 | Paper 10.25 describes the tools for reviewing the T10 master receipt. These tools expose receipt binding, transport classification, local witness}\\
\noindent\texttt{| Additional evidence cards | 34 total | The complete card inventory is maintained in the archive evidence matrix. | Cards are source-discovery aids until bound to paper-local evidence. |}\\

\subsection*{Journal Apparatus}

\subsubsection*{Article Type And Intended Venue Fit}

This manuscript is written as a formal-methods and mathematical-systems paper
inside the series *Formalizing LCR, Unifying Standard Models*. Its intended
reader is a technical reviewer who needs claim boundaries, reproducibility
routes, and cross-paper dependencies to be visible without relying on private
context.

\subsubsection*{Structured Contribution Statement}

\noindent\texttt{| Item | Statement |}\\
\noindent\texttt{|---|---|}\\
\noindent\texttt{| Publication ID | `FLCR-02` |}\\
\noindent\texttt{| Legacy source slot(s) | `01` |}\\
\noindent\texttt{| Ribbon role | order-1 slot-1: start the active one-dimensional action / enumerate the center state |}\\
\noindent\texttt{| Proof form | enumeration proof and identity preservation |}\\
\noindent\texttt{| Received state | state emitted by prior slot 00 and reopened at original slot 01 (LCR Chain Carrier) |}\\
\noindent\texttt{| Produced state | formal result of original slot 01 plus the same-family lift path toward slot 11 |}\\

\subsubsection*{Claim-Evidence Table}

\noindent\texttt{| Claim | Lane | Evidence to attach | Boundary |}\\
\noindent\texttt{|---|---|---|---|}\\
\noindent\texttt{| Theorem 2.1 | `receipt\_bound\_internal\_result` | Receipt, source card, validator, citation, or CAM route named in the paper manifest | The LCR carrier is the minimal address-preserving carrier for one center and two dis}\\
\noindent\texttt{| Theorem 2.2 | `receipt\_bound\_internal\_result` | Receipt, source card, validator, citation, or CAM route named in the paper manifest | In the binary chart, reversal is an involution and partitions the carrier into fixed}\\
\noindent\texttt{| Proposition 2.3 | `standard\_theorem\_citation\_bound\_result` | Receipt, source card, validator, citation, or CAM route named in the paper manifest | The shell profile of \{0,1\}\textasciicircum{}3 is the binomial profile 1,3,3,1. |}\\

\subsubsection*{Figure Plan}

\noindent\texttt{| Figure | Purpose | Caption |}\\
\noindent\texttt{|---|---|---|}\\
\noindent\texttt{| FLCR-02-F1 | State enumeration chart | Schematic showing how `FLCR-02` turns its received state into the produced state while preserving claim lanes and residue boundaries. |}\\
\noindent\texttt{| FLCR-02-F2 | Evidence routing map | Diagram of source papers, archive cards, CAM/crystal routes, validators, and workbook replay paths used by this manuscript. |}\\
\noindent\texttt{| FLCR-02-F3 | Same-family lift context | Diagram placing this paper in the original `00-79` ribbon family and showing predecessor/successor slots. |}\\

\subsubsection*{In-Text Figure FLCR-02-F1: State enumeration chart}

\begin{verbatim}
received state
  -> State enumeration chart
  -> formal claim lane
  -> evidence object / receipt / citation
  -> produced state
  -> remaining residue
\end{verbatim}

\subsubsection*{Table Plan}

\noindent\texttt{| Table | Purpose |}\\
\noindent\texttt{|---|---|}\\
\noindent\texttt{| FLCR-02-T1 | State inventory table |}\\
\noindent\texttt{| FLCR-02-T2 | Claim-lane/evidence-boundary table |}\\
\noindent\texttt{| FLCR-02-T3 | Archive evidence card and source-backed upgrade table |}\\
\noindent\texttt{| FLCR-02-T4 | Workbook replay and falsifier table |}\\

\subsubsection*{Worked Example}

Enumerate the local carrier or admission candidates and classify each row before interpretation. The example must name the input object, the operation,
the evidence object, the allowed revised claim, and the remaining boundary.

\subsubsection*{Nomenclature And Glossary}

\noindent\texttt{| Term | Paper-local meaning |}\\
\noindent\texttt{|---|---|}\\
\noindent\texttt{| CAM | Defined by this paper as part of the `enumeration\_action` proof role; final definition is recorded in the paper-local glossary. |}\\
\noindent\texttt{| addressable slot | Defined by this paper as part of the `enumeration\_action` proof role; final definition is recorded in the paper-local glossary. |}\\
\noindent\texttt{| center state | Defined by this paper as part of the `enumeration\_action` proof role; final definition is recorded in the paper-local glossary. |}\\
\noindent\texttt{| claim lane | Defined by this paper as part of the `enumeration\_action` proof role; final definition is recorded in the paper-local glossary. |}\\
\noindent\texttt{| crystal | Defined by this paper as part of the `enumeration\_action` proof role; final definition is recorded in the paper-local glossary. |}\\
\noindent\texttt{| enumeration | Defined by this paper as part of the `enumeration\_action` proof role; final definition is recorded in the paper-local glossary. |}\\
\noindent\texttt{| finite domain | Defined by this paper as part of the `enumeration\_action` proof role; final definition is recorded in the paper-local glossary. |}\\
\noindent\texttt{| receipt | Defined by this paper as part of the `enumeration\_action` proof role; final definition is recorded in the paper-local glossary. |}\\

\subsubsection*{Appendix FLCR-02-A: Evidence Package}

The appendix records:

1. Legacy source excerpts or source-card references used in the final claim ledger.
2. Receipt and verifier paths, including pass/fail/partial status.
3. CAM/crystal query or projection rows used by the paper, with source identity and boundary.
4. Standards-window and NSIT evidence used for conformance or routing.
5. Falsifier, calibration, or external-data rows that remain unresolved.

\subsubsection*{Data And Code Availability}

The publication package contains the canonical Markdown sources, manifests, source-routing matrices, validation reports, and generated PDF/TeX outputs.

Code and verifier references are cited through manifests, archive evidence cards, receipt catalogs, or workbook replay tables. External datasets or
standards citations must be attached before any calibration-bound claim is
promoted.

\subsubsection*{Reviewer Checklist}

\noindent\texttt{| Check | Reviewer question |}\\
\noindent\texttt{|---|---|}\\
\noindent\texttt{| Claim lane | Does every strong claim declare its lane and evidence type? |}\\
\noindent\texttt{| Boundary | Does the paper preserve what it does not prove? |}\\
\noindent\texttt{| Reproducibility | Is there a receipt, verifier, source card, citation, or replay table for the claim? |}\\
\noindent\texttt{| Cross-paper import | Does the paper cite the prior FLCR/OPR slot before consuming its result? |}\\
\noindent\texttt{| Interpretation | Does the analog layer preserve the addressed state while changing labels/access paths? |}\\

\subsubsection*{Declarations}

No human-subjects, clinical, or animal-research protocol is asserted by this
paper. External empirical claims require separate dataset, calibration, or
domain-validation attachments. Competing-interest and funding statements are recorded in the submission metadata.

\subsection*{11. Publication Integration Addendum}

\subsubsection*{11.1 Role In The Series}

`FLCR-02` (`The LCR Carrier`) occupies the `enumeration\_action` role at lift depth `0`.
The paper receives state emitted by prior slot 00 and reopened at original slot 01 (LCR Chain Carrier). It produces formal result of original slot 01 plus the same-family lift path toward slot 11. The state transition is:
select the active center and enumerate the addressable state space at the current lift.

The paper-local proof form is:

\begin{verbatim}
enumeration proof and identity preservation
\end{verbatim}

The integration result is an identity-preserving inventory of admissible state objects. This addendum records how that result is
consumed by the publication series without relying on editorial instructions or
private working context.

\subsubsection*{11.2 Formal Carrier}

\noindent\texttt{| Formal carrier | Function |}\\
\noindent\texttt{|---|---|}\\
\noindent\texttt{| finite set enumeration | Formal carrier for the paper-local result. |}\\
\noindent\texttt{| ordered tuples and projections | Formal carrier for the paper-local result. |}\\
\noindent\texttt{| group actions on finite carriers | Formal carrier for the paper-local result. |}\\
\noindent\texttt{| minimality or lower-bound arguments | Formal carrier for the paper-local result. |}\\

The LCR, CAM, crystal, forge, and analog vocabulary functions as a typed access
layer over these carriers. A relabeling is admissible only when the addressed
object, evidence lane, boundary, and downstream import rule remain attached.

\subsubsection*{11.3 Claim Ledger}

\noindent\texttt{| Claim | Lane | Statement |}\\
\noindent\texttt{|---|---|---|}\\
\noindent\texttt{| Theorem 2.1 | `receipt\_bound\_internal\_result` | The LCR carrier is the minimal address-preserving carrier for one center and two distinguishable boundary slots. |}\\
\noindent\texttt{| Theorem 2.2 | `receipt\_bound\_internal\_result` | In the binary chart, reversal is an involution and partitions the carrier into fixed states and reversal pairs. |}\\
\noindent\texttt{| Proposition 2.3 | `standard\_theorem\_citation\_bound\_result` | The shell profile of \{0,1\}\textasciicircum{}3 is the binomial profile 1,3,3,1. |}\\

These claims are consumed at their stated lane. A stronger downstream statement
creates a new claim envelope with its own evidence object.

\subsubsection*{11.4 Evidence Package}

\noindent\texttt{| Evidence class | Routed items | Status |}\\
\noindent\texttt{|---|---|---|}\\
\noindent\texttt{| Legacy sources | `01\_LCR\_Chain\_Carrier.md`, `supplements/01\_INTERNAL\_REASONING\_SUPPLEMENT.md` | routed evidence |}\\
\noindent\texttt{| Evidence inputs | `CAM\_CLAIM\_100\_SCORE\_AUDIT.md`, `NSIT\_TOOL\_CLOSURE\_MATRIX.md`, `NSIT\_REASONED\_CLOSURE\_PASS.md`, `PAPER\_REWORKS\_CRYSTAL\_PROJECTION.json` | routed evidence |}\\
\noindent\texttt{| CAM/crystal sources | `PAPER\_REWORKS\_CRYSTAL\_PROJECTION.json`, `AGENT\_CRYSTAL\_WORK\_HARVEST.json`, `runtime standards artifact`, `notebook-derived source bundle` | routed evidence |}\\
\noindent\texttt{| Standards-window inputs | `window\_summary.json`, `node DBs`, `window DBs` | routed evidence |}\\
\noindent\texttt{| Receipts | external requirement | not attached in this source package |}\\
\noindent\texttt{| Validators | external requirement | not attached in this source package |}\\

Standards-window, runtime, and notebook-derived artifacts are treated
as evidence-routing surfaces. They support source discovery, conformance
checking, and reapplication, but they do not replace citations, receipts,
normal-form derivations, calibration rows, or falsifiers.

\subsubsection*{11.5 Results Representation}

The paper's results section is organized around five publication objects:

1. the exact object acted on at this slot;
2. the admissible operation or transformation;
3. the receipt, enumeration, derivation, lookup, or external theorem supporting
   the operation;
4. the residue or boundary left after the result;
5. the downstream import rule.

\subsubsection*{11.6 Figures And Tables}

\noindent\texttt{| Figure | Role |}\\
\noindent\texttt{|---|---|}\\
\noindent\texttt{| FLCR-02-F1: State enumeration chart | Visualizes the proof object, routing, or boundary. |}\\
\noindent\texttt{| FLCR-02-F2: Evidence routing map | Visualizes the proof object, routing, or boundary. |}\\
\noindent\texttt{| FLCR-02-F3: Same-family lift context | Visualizes the proof object, routing, or boundary. |}\\

\noindent\texttt{| Table | Role |}\\
\noindent\texttt{|---|---|}\\
\noindent\texttt{| FLCR-02-T1: State inventory table | Records claim lanes, evidence, sources, or falsifiers. |}\\
\noindent\texttt{| FLCR-02-T2: Claim-lane/evidence-boundary table | Records claim lanes, evidence, sources, or falsifiers. |}\\
\noindent\texttt{| FLCR-02-T3: Archive evidence card and source-backed upgrade table | Records claim lanes, evidence, sources, or falsifiers. |}\\
\noindent\texttt{| FLCR-02-T4: Workbook replay and falsifier table | Records claim lanes, evidence, sources, or falsifiers. |}\\

\subsubsection*{11.7 Limits And Falsifiers}

The paper proves only the object and domain declared in its claim ledger.
Physical, Standard Model, observer, calibration, or synthesis claims require
the additional lane evidence stated by their destination papers. CAM/crystal
and forge language is operational only when represented as an address, lookup,
receipt, validator, or reapplication path.
\end{document}
